\documentclass[11pt]{article}

\usepackage[margin=1in]{geometry}
\usepackage{amsmath, amssymb, amsthm, mathtools}
\usepackage{booktabs}
\usepackage{enumitem}
\usepackage{hyperref}
\hypersetup{colorlinks=true, linkcolor=blue, citecolor=blue, urlcolor=blue}

% ---------------------------------------------------------------------------
% Notation, matching https://slinderman.github.io/stats305c/em/ and /hmms/
% ---------------------------------------------------------------------------
\newcommand{\N}{\mathcal{N}}
\newcommand{\Cat}{\mathrm{Cat}}
\newcommand{\KL}{D_{\mathrm{KL}}}
\newcommand{\E}{\mathbb{E}}
\DeclareMathOperator*{\argmax}{arg\,max}

\newtheorem*{lemma*}{Lemma}

\title{EM and the Forward--Backward Algorithm for Hidden Markov Models}
\author{Scott Linderman \\ \small Methods in Computational Neuroscience, MBL --- 2026}
\date{}

\begin{document}
\maketitle

\begin{abstract}
\noindent
These notes derive the forward-backward algorithm and EM (Baum-Welch) for
hidden Markov models directly, by maximum likelihood. This is the machinery
underneath every \texttt{.fit(method="em")} call in this afternoon's coding
tutorial --- today we derive what that call is actually doing, down to the
weighted linear regression at the heart of the tutorial's AR-HMM M-step.
Based on
\href{https://slinderman.github.io/stats305c/em/}{\texttt{stats305c/em}} and
\href{https://slinderman.github.io/stats305c/hmms/}{\texttt{stats305c/hmms}}.
\end{abstract}

\tableofcontents

% =============================================================================
\section{Hidden Markov Models}
% =============================================================================

A \textbf{hidden Markov model (HMM)} pairs a Markov chain of discrete latent
states $z_1, \ldots, z_T \in \{1,\ldots,K\}$ with a sequence of observations
$x_1, \ldots, x_T$, one emitted per state. The joint distribution
factors as
\begin{equation}
p(z_{1:T}, x_{1:T}; \Theta)
= p(z_1; \pi_0) \prod_{t=2}^T p(z_t \mid z_{t-1}; P) \prod_{t=1}^T p(x_t \mid \theta_{z_t}),
\label{eq:hmm-joint}
\end{equation}
with three parameter groups $\Theta = (\pi_0, P, \{\theta_k\}_{k=1}^K)$:
\begin{enumerate}[itemsep=2pt]
  \item \textbf{Initial distribution:} $z_1 \sim \Cat(\pi_0)$.
  \item \textbf{Transition matrix:} $P \in [0,1]^{K\times K}$
    (row-stochastic), $P_{ij} = p(z_{t+1}{=}j \mid z_t{=}i)$.
  \item \textbf{Emissions:} $x_t \sim p(\cdot \mid \theta_{z_t})$ ---
    e.g. Gaussian, Poisson, or, for this afternoon's tutorial, an
    \emph{autoregressive} Gaussian, $x_t \mid z_t{=}k, x_{t-1} \sim
    \N(A_k x_{t-1} + b_k,\, Q_k)$.
\end{enumerate}
We want the \textbf{maximum likelihood estimate} of $\Theta$: the parameters
that make the observed data $x_{1:T}$ most probable, marginalizing over
the states we never observe,
\begin{equation}
\hat\Theta = \argmax_\Theta \ \log p(x_{1:T}; \Theta),
\qquad
p(x_{1:T}; \Theta) = \sum_{z_{1:T}} p(z_{1:T}, x_{1:T}; \Theta).
\label{eq:mle}
\end{equation}
The sum in \eqref{eq:mle} is over all $K^T$ possible state sequences ---
already intractable to evaluate for, say, $K=6$ and a few hundred time
steps, let alone to differentiate for gradient ascent. \textbf{EM} sidesteps
this by exploiting the Markov structure of \eqref{eq:hmm-joint}: it turns
out we only ever need two low-dimensional summaries of the posterior over
paths, not the whole $K^T$-entry table, and those summaries can be computed
by the \textbf{forward-backward algorithm} in $O(TK^2)$ time. That's the
plan for the rest of these notes: derive EM in general, including the
M-step for a simple Gaussian emission model (Section~2), then derive
forward-backward as the way to compute the summaries EM's E-step needs
(Section~3). The M-step for the tutorial's autoregressive emissions is left
as an exercise at the end of Section~2.


% =============================================================================
\section{The EM Algorithm}
% =============================================================================

\subsection{Jensen's inequality and the ELBO}

\begin{lemma*}[Jensen's inequality]
For any concave function $f$ and random variable $Y$,
$f(\E[Y]) \geq \E[f(Y)]$, with equality iff $Y$ is constant (or $f$ is
linear). Since $\log$ is concave, $\log \E[g] \geq \E[\log g]$ for any
positive random variable $g$.
\end{lemma*}

Introduce an \textbf{auxiliary distribution} $q(z_{1:T})$ over the whole
hidden path --- for now, completely arbitrary, with the same support as the
true posterior. Multiply and divide by $q$ inside the sum in
\eqref{eq:mle}, then apply Jensen's inequality:
\begin{align}
\log p(x_{1:T}; \Theta) 
&= \log \sum_{z_{1:T}} p(z_{1:T}, x_{1:T}; \Theta) \nonumber \\
&= \log \sum_{z_{1:T}} q(z_{1:T}) \, \frac{p(z_{1:T}, x_{1:T}; \Theta)}{q(z_{1:T})} \nonumber  \\
&= \log \E_q\!\left[\frac{p(z_{1:T}, x_{1:T}; \Theta)}{q(z_{1:T})}\right] \notag \\
&\geq \E_q\Bigl[\log p(z_{1:T}, x_{1:T}; \Theta) - \log q(z_{1:T})\Bigr] \nonumber \\
&\triangleq \mathcal{L}(\Theta, q).
\label{eq:elbo-def}
\end{align}
$\mathcal{L}(\Theta,q)$ is the \textbf{evidence lower bound (ELBO)}. It holds
for \emph{any} $q$, so we've replaced one hard problem, maximizing
$\log p(x_{1:T};\Theta)$ over $\Theta$, with an easier joint problem:
maximize $\mathcal{L}(\Theta,q)$ over $(\Theta, q)$ by coordinate ascent.
That coordinate ascent is EM --- E-step optimizes $q$, M-step optimizes
$\Theta$ --- and it is preferable to direct gradient ascent on
\eqref{eq:mle} because, as we'll see, each coordinate update has a simple,
closed form.

\subsection{The E-step}

Note that $\log p(z_{1:T}, x_{1:T}; \Theta) = \log p(z_{1:T} \mid x_{1:T}; \Theta) + \log p(x_{1:T}; \Theta)$. Expanding the joint distribution inside the
expectation in \eqref{eq:elbo-def} and rearranging terms --- a procedure we might call \textit{ELBO surgery} --- we obtain another form of the bound,
\begin{equation}
\mathcal{L}(\Theta, q) = \log p(x_{1:T}; \Theta)
- \KL\bigl(q(z_{1:T}) \,\|\, p(z_{1:T} \mid x_{1:T}; \Theta)\bigr),
\label{eq:elbo-kl}
\end{equation}
where $\KL(q\|p) = \sum_z q(z)\log\frac{q(z)}{p(z)}$ is the \textbf{Kullback-Leibler (KL) divergence}. Since the KL divergence is non-negative and zero iff $q=p$,
we see that $\mathcal{L} \leq \log p(x_{1:T};\Theta)$ always, and
maximizing $\mathcal L$ over $q$ (with $\Theta$ fixed) means minimizing that
KL divergence. The unique minimizer is
\begin{equation}
\boxed{q^\star(z_{1:T}) = p(z_{1:T} \mid x_{1:T}; \Theta),}
\label{eq:e-step}
\end{equation}
the exact posterior over the whole path, at which point the bound is tight:
$\mathcal L(\Theta, q^\star) = \log p(x_{1:T};\Theta)$. This is the
\textbf{E-step}.

\textbf{The catch:} $q^\star$ is a distribution over $K^T$ paths --- we
can't just write down a table of that size, let alone use it directly in
the M-step below. Fortunately, as we will see, the M-step only needs \textit{posterior marginals}, and we can compute those efficiently thanks to the Markovian nature of the model.

\subsection{The M-step}

The M-step maximizes the ELBO with respect to the parameters $\Theta$, holding the posterior $q^\star$ fixed. Let's focus on the emission parameters $\{\theta_k\}_{k=1}^K$. As
a function of those parameters, the ELBO is
\begin{align}
  \mathcal{L}(\Theta, q)
  &= \E_q\left[ \log p(x_{1:T}, z_{1:T}; \Theta)\right] + \mathrm{c} \\
  &= \sum_{t=1}^T \sum_k \underbrace{q(z_t = k)}_{\triangleq \gamma_t(k)} \log p(x_t ; \theta_k) + \mathrm{c},
\end{align}
where $\gamma_t(k)$ is the probability of being in state $k$ at time $t$ under the posterior $q$.
Maximizing with respect to the parameters $\theta_k$ is simply a weighted maximum likelihood estimation problem. Often, the solution can be found in closed form.

\paragraph{Example: Gaussian emissions}

Let's consider a simple case where the emission model is Gaussian with mean $\theta_k$,
\begin{align}
  p(x_t; \theta_k) &= \N(x_t \mid \theta_k, \sigma^2 I).
\end{align}
The ELBO as a function of $\theta_k$ is
\begin{align}
  \mathcal{L}(\Theta, q) &=
  \sum_{t=1}^T \gamma_t(k) \log \N(x_t \mid \theta_k, \sigma^2 I) \\
  &= \sum_{t=1}^T \gamma_t(k) \left( -\frac{1}{2\sigma^2} (x_t - \theta_k)^\top (x_t - \theta_k)\right).
\end{align}
It's a little bit of calculus to show that this objective is maximized at
\begin{align}
  \hat\theta_k &= \frac{\sum_{t=1}^T \gamma_t(k) x_t}{\sum_{t=1}^T \gamma_t(k)}.
\end{align}
Intuitively, we simply set the mean parameter equal to the average of observations, weighted by the probability of being in state $k$.

\paragraph{Exercise}
Now consider an autoregressive observation model (dropping the bias term for simplicity; see the remark below), $p(x_t \mid x_{t-1}; \theta_k) = \N(x_t \mid A_k x_{t-1}, Q_k)$.
Show that the optimal dynamics parameters are obtained by the
\textbf{weighted normal equations}, with solution
\begin{equation}
\boxed{\hat A_k
= \left(\sum_{t=1}^T \gamma_t(k)\, x_t x_{t-1}^\top\right)
  \left(\sum_{t=1}^T \gamma_t(k)\, x_{t-1}x_{t-1}^\top\right)^{-1}.}
\label{eq:mstep-A}
\end{equation}
This is exactly \textbf{weighted least squares}: regress $x_t$ on
$x_{t-1}$, weighting time step $t$ by how confident the
E-step is that $z_t = k$. (The tutorial's actual fit includes a bias term
too, $A_kx_{t-1}+b_k$; recover it with the usual trick for affine
regression --- augment $x_{t-1}$ with a constant $1$ and fold $b_k$ into
$A_k$ as an extra column.)


\subsection{EM as a minorize-maximize algorithm}

Each iteration: 
\begin{enumerate}
  \item \textbf{E-step} (minorize) --- build
$\mathcal L(\Theta^{(i)}, q)$, which touches $\log p(x_{1:T};\Theta)$ at
$\Theta = \Theta^{(i)}$; 
  \item \textbf{M-step} (maximize) --- set
$\Theta^{(i+1)} = \argmax_\Theta \mathcal L(\Theta, q^{(i)})$. 
\end{enumerate}
Since the M-step can only increase $\mathcal L$, and $\mathcal L \leq \log
p(x_{1:T};\Theta)$ always, the marginal log-likelihood is
\textbf{guaranteed to increase monotonically} every iteration, with no step
size to tune. EM converges to a \emph{local} optimum, and (like $K$-means)
is sensitive to initialization; random restarts are standard practice.


% =============================================================================
\section{The Forward-Backward Algorithm}
% =============================================================================

Our target quantity is the marginal $\gamma_t(k) = p(z_t{=}k \mid
x_{1:T})$ needed for the M-step above. The key to computing it
efficiently is the following factorization.

\textbf{Claim:}
\begin{equation}
p(z_t{=}k \mid x_{1:T}) \;\propto\;
\underbrace{p(z_t{=}k \mid x_{1:t-1})}_{\triangleq\ \alpha_t(k)}
\;\; p(x_t \mid z_t{=}k) \;\;
\underbrace{p(x_{t+1:T} \mid z_t{=}k)}_{\triangleq\ \beta_t(k)}.
\label{eq:gamma-factorization}
\end{equation}

\begin{proof}[Proof, by two applications of Bayes' rule]
Split $x_{1:T} = (x_{1:t-1}, x_t, x_{t+1:T})$
and apply Bayes' rule twice, each time treating one more block of
observations as newly-arrived evidence. Conditioning on $x_{t+1:T}$
last,
\[
p(z_t{=}k \mid x_{1:t-1}, x_t, x_{t+1:T})
\;\propto\; p(x_{t+1:T} \mid z_t{=}k, x_{1:t-1}, x_t) \;
p(z_t{=}k \mid x_{1:t-1}, x_t).
\]
In an HMM, once we know $z_t$, everything that happens afterward is
independent of everything that happened before --- the defining Markov
property --- so $p(x_{t+1:T} \mid z_t{=}k, x_{1:t-1}, x_t) =
p(x_{t+1:T}\mid z_t{=}k) = \beta_t(k)$. Applying Bayes' rule once more to
the remaining factor, now splitting off $x_t$,
\[
p(z_t{=}k \mid x_{1:t-1}, x_t) \;\propto\;
p(x_t \mid z_t{=}k, x_{1:t-1}) \; p(z_t{=}k \mid x_{1:t-1}).
\]
Since $x_t$ depends only on $z_t$, $p(x_t \mid z_t{=}k, x_{1:t-1}) =
p(x_t\mid z_t{=}k)$, and $p(z_t{=}k\mid x_{1:t-1}) = \alpha_t(k)$ by
definition. Multiplying the pieces back together gives
\eqref{eq:gamma-factorization}.
\end{proof}

So the whole problem reduces to computing two quantities recursively: the
\textbf{forward message} $\alpha_t(k) \triangleq p(z_t{=}k\mid x_{1:t-1})$
and the \textbf{backward message} $\beta_t(k) \triangleq p(x_{t+1:T}\mid
z_t{=}k)$.

\subsection{The forward algorithm}

\textbf{Claim:} $\displaystyle \alpha_{t+1}(j) = \frac{1}{C_t}\sum_{k=1}^K
P_{kj}\, \alpha_t(k)\, p(x_t\mid z_t{=}k)$, where $C_t \triangleq
p(x_t\mid x_{1:t-1}) = \sum_k \alpha_t(k)\, p(x_t\mid z_t{=}k)$.

\begin{proof}[Proof, by induction on $t$]
\textit{Base case:} $\alpha_1(k) = p(z_1{=}k\mid x_{1:0}) = p(z_1{=}k) =
\pi_{0,k}$ --- there is no data to condition on before $t=1$.

\textit{Inductive step:} suppose $\alpha_t(k)$ is known for every $k$.
First, incorporate the new observation $x_t$ via Bayes' rule to get the
\emph{filtered} distribution,
\[
p(z_t{=}k \mid x_{1:t}) = \frac{p(x_t\mid z_t{=}k)\,\alpha_t(k)}
{\sum_{k'} p(x_t\mid z_t{=}k')\,\alpha_t(k')}
= \frac{p(x_t\mid z_t{=}k)\,\alpha_t(k)}{C_t}.
\]
Then predict one step ahead: by the law of total probability and the
Markov property ($z_{t+1}$ depends on the past only through $z_t$),
\[
\alpha_{t+1}(j) = p(z_{t+1}{=}j\mid x_{1:t})
= \sum_k p(z_{t+1}{=}j\mid z_t{=}k)\, p(z_t{=}k\mid x_{1:t})
= \sum_k P_{kj}\, \frac{\alpha_t(k)\, p(x_t\mid z_t{=}k)}{C_t}. \qedhere
\]
\end{proof}

Writing $\ell_t = [p(x_t\mid z_t{=}1),\ldots,p(x_t\mid z_t{=}K)]^\top$
for the vector of emission likelihoods at time $t$, this is the matrix-vector
recursion
\begin{equation}
\alpha_{t+1} = \frac{P^\top(\alpha_t \odot \ell_t)}{C_t},
\qquad C_t = \alpha_{t}^\top \ell_{t},
\label{eq:forward-recursion}
\end{equation}
costing $O(K^2)$ per step ($O(TK^2)$ total, versus $O(K^T)$ for naive
enumeration). Notice normalization wasn't an afterthought bolted on for
numerical stability --- it fell straight out of the derivation, because
$\alpha_t$ is a genuine conditional distribution (Bayes' rule) rather than
an unnormalized joint probability. As a bonus, $C_t = p(x_t\mid
x_{1:t-1})$ is exactly the one-step-ahead predictive likelihood, so by
the chain rule we get the \textbf{marginal log-likelihood for free}:
\begin{equation}
\log p(x_{1:T}) = \sum_{t=1}^T \log C_t.
\label{eq:marginal-lik-hmm}
\end{equation}
This is the quantity we track as ``the log-likelihood'' during EM.

\subsection{The backward algorithm}

\textbf{Claim:} $\displaystyle \beta_t(k) = \sum_{j=1}^K P_{kj}\,
p(x_{t+1}\mid z_{t+1}{=}j)\, \beta_{t+1}(j)$.

\begin{proof}[Proof, by induction on $t$, downward from $T$]
\textit{Base case:} $\beta_T(k) = p(x_{T+1:T}\mid z_T{=}k) = 1$, since
the ``future after $T$'' is the empty sequence of observations, which is
certain.

\textit{Inductive step:} suppose $\beta_{t+1}(j)$ is known for every $j$.
Introduce $z_{t+1}$ by marginalizing,
\[
\beta_t(k) = p(x_{t+1}, x_{t+2:T}\mid z_t{=}k)
= \sum_j p(z_{t+1}{=}j\mid z_t{=}k)\, p(x_{t+1}\mid z_{t+1}{=}j, z_t{=}k)\,
p(x_{t+2:T}\mid z_{t+1}{=}j, z_t{=}k, x_{t+1}).
\]
By the Markov property, both $x_{t+1}$ and $x_{t+2:T}$ depend on $z_t$
only through $z_{t+1}$, so this reduces to $\sum_j P_{kj}\,
p(x_{t+1}\mid z_{t+1}{=}j)\, \beta_{t+1}(j)$. \qedhere
\end{proof}

In matrix form, $\beta_t = P(\beta_{t+1}\odot\ell_{t+1})$. The
backward pass runs from $t=T$ down to $t=1$, the mirror image of the
forward pass.

\textbf{A note on normalization.} Unlike $\alpha_t$, $\beta_t$ is not a
normalized distribution over $k$ --- it's a product of $T-t$ likelihood
terms, so it can just as easily underflow for large $T$. But we're free to
rescale it: since $\gamma_t(k) \propto \alpha_t(k)\,\ell_{t,k}\,\beta_t(k)$
only needs $\beta_t(k)$ up to an overall constant, one that doesn't depend
on $k$, we can renormalize $\beta_t(\cdot)$ by any convenient constant
$B_t$ at every step of the recursion --- e.g.\ $\beta_t(k) \leftarrow
\beta_t(k)/B_t$ with $B_t = \sum_k \beta_t(k)$ --- without changing
$\gamma_t(k)$ once it's normalized over $k$. This is the same numerical
trick as for $\alpha_t$, just applied by hand here instead of falling out
of the derivation automatically.

\subsection{Putting it together}

One forward pass and one backward pass give $\alpha_t(k)$ and $\beta_t(k)$
for every $t$, and hence, via \eqref{eq:gamma-factorization}, the posterior
marginal
\begin{equation}
\gamma_t(k) = p(z_t{=}k \mid x_{1:T}) \;\propto\; \alpha_t(k)\, \ell_{t,k}\, \beta_t(k).
\label{eq:posterior-marginal}
\end{equation}
This is the \textbf{E-step}, computed in $O(TK^2)$ time, without ever
writing down the $K^T$-entry posterior over paths.

\paragraph{Exercise} The transition-matrix M-step needs one more quantity:
the pairwise marginal $\xi_t(i,j) \triangleq p(z_t{=}i, z_{t+1}{=}j \mid
x_{1:T})$. Using the same conditional-independence argument as the proof
of \eqref{eq:gamma-factorization} --- but now keeping \emph{both} $z_t$ and
$z_{t+1}$ fixed, rather than marginalizing the transition between them ---
show that
\[
\xi_t(i,j) \;\propto\; \alpha_t(i)\, \ell_{t,i}\, P_{ij}\, \ell_{t+1,j}\, \beta_{t+1}(j)
\]
(normalize over $i,j$ to get an honest probability). This is exactly the
weight $\hat P_{ij} \propto \sum_{t=1}^{T-1} \xi_t(i,j)$ needs in the
transition-matrix M-step.


\begin{thebibliography}{9}
\bibitem{baum1970maximization}
L.~E. Baum, T.~Petrie, G.~Soules, and N.~Weiss.
A maximization technique occurring in the statistical analysis of
probabilistic functions of Markov chains.
\emph{The Annals of Mathematical Statistics}, 41(1):164--171, 1970.

\bibitem{rabiner1986introduction}
L.~R. Rabiner and B.~H. Juang.
An introduction to hidden Markov models.
\emph{IEEE ASSP Magazine}, 3(1):4--16, 1986.

\bibitem{wiltschko2015mapping}
A.~B. Wiltschko, M.~J. Johnson, G.~Iurilli, R.~E. Peterson, J.~M. Katon,
S.~L. Pashkovski, V.~E. Abraira, R.~P. Adams, and S.~R. Datta.
Mapping sub-second structure in mouse behavior.
\emph{Neuron}, 88(6):1121--1135, 2015.
\end{thebibliography}

\end{document}
